\documentclass[14pt]{extarticle}
\usepackage{fontspec}
\usepackage{polyglossia}
\usepackage{amsmath}
\usepackage{amssymb}
\usepackage{xcolor}
\usepackage{fancyhdr}
\usepackage{graphicx}
\usepackage{listings}
\usepackage[most]{tcolorbox}
% \usepackage{setspace}
\usepackage{pifont}
\usepackage{enumitem}
\usepackage{titlesec}
\usepackage[bottom]{footmisc}
\usepackage{titling}
\usepackage{minted}
\usepackage{etoolbox}
\usepackage{array}
\usepackage{longtable} % page-breakable tables (must load before polyglossia/bidi RTL setup)
\usepackage{extsizes}
\usepackage{float}
\usepackage{booktabs}
\usepackage{hyperref}
\usepackage[onehalfspacing]{setspace}
\newif\ifwhiteblack
\ifdefined\setwhiteblack
  \setwhiteblack
\fi

\ifwhiteblack
  \usepackage[a4paper, landscape, top=1.25cm, bottom=1.25cm, left=1.25cm, right=1.25cm, headheight=1.5cm, headsep=0.5cm]{geometry}
\else
  \usepackage[a4paper, top=2.7cm, bottom=2.7cm, left=2.5cm, right=2.5cm, headheight=1.5cm, headsep=0.5cm]{geometry}
\fi
\usepackage{tikz}
\usepackage{datetime2}

\usetikzlibrary{automata, positioning, arrows.meta, arrows}
\providecommand{\LTcaptype}{table}

\newfontfamily\emoji{Segoe UI Emoji}

\pagestyle{fancy}

\setmainlanguage[numerals=western]{arabic}
\setotherlanguage{english}
\newfontfamily\arabicfont[Script=Arabic]{Calibri}
\newfontfamily\arabicfonttt[Script=Arabic]{Courier New}



\ifwhiteblack
\lstset{
  language=[Sharp]C,
  numbers=left,
  stepnumber=1,
  numbersep=8pt,
  frame=single,
  basicstyle=\ttfamily\small,
  keywordstyle=\bfseries,
  stringstyle=\ttfamily,
  commentstyle=\itshape
}
\else
\lstset{
  language=[Sharp]C,
  numbers=left,
  stepnumber=1,
  numbersep=8pt,
  frame=single,
  basicstyle=\ttfamily\small,
  keywordstyle=\color{blue},
  stringstyle=\color{red},
  commentstyle=\color{green!50!black}
}
\fi

\newif\ifdetailed
\ifdefined\setdetailed
  \setdetailed
\fi

\newif\ifwithsols
\ifdefined\setwithsols
  \setwithsols
\fi

\ifwhiteblack
  \AtBeginDocument{
    \LARGE
    \usemintedstyle{bw}
  }
\fi

% Default minted look: compact code that wraps long lines.
% Not applied in black-and-white landscape mode, which already uses \LARGE.
\ifwhiteblack\else
  \setminted{fontsize=\small, breaklines=true}
\fi

% unified tcolorboxes for articles
\ifwhiteblack
\tcbset{colback=white, colframe=black, fonttitle=\bfseries, boxrule=1pt}
\newtcolorbox{boxDef}[1][]{title={{\emoji📘} تعريف\ifx\\#1\\\else ~#1\fi :}}
\newtcolorbox{boxExercise}[1][]{title={{\emoji🧩} تمرين\ifx\\#1\\\else ~#1\fi :}}
\newtcolorbox{boxExample}[1][]{title={{\emoji📝} مثال\ifx\\#1\\\else ~#1\fi :}}
\newtcolorbox{boxNote}[1][]{title={{\emoji✨} ملاحظة\ifx\\#1\\\else ~#1\fi :}}
\newtcolorbox{boxAttention}[1][]{title={{\emoji🔔} تنبيه\ifx\\#1\\\else ~#1\fi :}}
\newtcolorbox{boxWarning}[1][]{title={{\emoji⚡} ملاحظة هامة\ifx\\#1\\\else ~#1\fi :}}
\newtcolorbox{boxSolution}[1][]{title={{\emoji✅} حل\ifx\\#1\\\else ~#1\fi :}}
\newtcolorbox{boxSymbol}[1][]{title={{\emoji🔣} رمز\ifx\\#1\\\else ~#1\fi :}}
\newtcolorbox{boxHint}[1][]{title={{\emoji💡} تلميح\ifx\\#1\\\else ~#1\fi :}}
\else
\tcbset{colback=white, colframe=black, fonttitle=\bfseries, boxrule=0.8pt}
\newtcolorbox{boxDef}[1][]{colback=blue!5!white,colframe=blue!75!black,
  title={{\emoji📘} تعريف\ifx\\#1\\\else ~#1\fi :}}
\newtcolorbox{boxExercise}[1][]{colback=cyan!5!white,colframe=cyan!70!black,
  title={{\emoji🧩} تمرين\ifx\\#1\\\else ~#1\fi :}}
\newtcolorbox{boxExample}[1][]{colback=yellow!5!white,colframe=orange!90!black,
  title={{\emoji📝} مثال\ifx\\#1\\\else ~#1\fi :}}
\newtcolorbox{boxNote}[1][]{colback=gray!10!white,colframe=black,
  title={{\emoji✨} ملاحظة\ifx\\#1\\\else ~#1\fi :}}
\newtcolorbox{boxAttention}[1][]{colback=magenta!10!white,colframe=magenta!80!black,
  title={{\emoji🔔} تنبيه\ifx\\#1\\\else ~#1\fi :}}
\newtcolorbox{boxWarning}[1][]{colback=red!5!white,colframe=red!75!black,
  title={{\emoji⚡} ملاحظة هامة\ifx\\#1\\\else ~#1\fi :}}
\newtcolorbox{boxSolution}[1][]{colback=green!5!white,colframe=green!60!black,
  title={{\emoji✅} حل\ifx\\#1\\\else ~#1\fi :}}
\newtcolorbox{boxSymbol}[1][]{colback=purple!5!white,colframe=purple!70!black,
  title={{\emoji🔣} رمز\ifx\\#1\\\else ~#1\fi :}}
\newtcolorbox{boxHint}[1][]{colback=teal!5!white,colframe=teal!60!black,
  title={{\emoji💡} تلميح\ifx\\#1\\\else ~#1\fi :}}
\fi


\ifwhiteblack
\tcbset{simplecode/.style={ colback=white, colframe=black, boxrule=1pt, arc=2pt, left=4pt,right=4pt,top=4pt,bottom=4pt}}
\else
\tcbset{simplecode/.style={ colback=gray!5, colframe=black!50, boxrule=0.4pt, arc=2pt, left=4pt,right=4pt,top=4pt,bottom=4pt}}
\fi
% Compact code boxes: tighter line spacing (the preamble uses onehalfspacing) and small vertical gaps.
% Must come AFTER the simplecode definition above.
\tcbset{simplecode/.append style={fontupper=\setstretch{0.95}, top=4pt, bottom=4pt, before skip=6pt, after skip=6pt}}
\newenvironment{boxCode}{\begin{tcolorbox}[simplecode]}{\end{tcolorbox}}

% Code cell used inside tables (keeps code left-to-right inside RTL tables).
% Use inside a fixed-width p{...} column: \begin{pycell}\begin{minted}...\end{minted}\end{pycell}
\newtcolorbox{pycell}{frame empty,colback=white,boxsep=0pt,left=2pt,right=2pt,top=2pt,bottom=2pt,width=\linewidth}

\newcolumntype{C}[1]{>{\centering\arraybackslash}p{#1}}

\makeatletter
\DTMnewdatestyle{dotted}{%
  \renewcommand{\DTMdisplaydate}[4]{\two@digits{##3}.\two@digits{##2}.##1}%
  \renewcommand{\DTMDisplaydate}{\DTMdisplaydate}%
}
\makeatother
\DTMsetdatestyle{dotted}

% redefine spaces after titles
\makeatletter
\renewcommand{\@maketitle}{%
  \begin{center}
    {\huge \bfseries \@title \par}%
    \vskip 0.2em % space between title and author
    {\large \@author \par}%
    % \vskip 0.2em % space between author and date
    % {\normalsize \@date \par}%
    {\small \hfill \begin{english}\DTMtoday\end{english} \par}%
  \end{center}
}
\makeatother

\fancyhf{} % clear default
\fancypagestyle{plain}{
  \fancyhf{}
  \fancyhead[L]{مدرسة التسامح الشاملة}
  % \fancyhead[L]{\includegraphics[height=1cm]{../../../images/logoTasamoh.png}}
  \fancyhead[R]{الأستاذ محمود اغبارية}
  \fancyfoot[C]{\thepage}
}

\fancyhead[L]{مدرسة التسامح الشاملة}
\fancyhead[R]{الأستاذ محمود اغبارية}
\fancyfoot[C]{\thepage}

\setcounter{tocdepth}{3} % only section subsection and subsubsection in TOC

\newcommand{\importpdfpage}[2]{%
    \thispagestyle{empty}
    \newgeometry{left=0cm,right=0cm,top=0cm,bottom=0cm}
    \includegraphics[page=#2,width=0.95\paperwidth,height=0.95\paperheight,keepaspectratio]{#1}
    \restoregeometry
}

\newcommand{\insertFullImg}[1]{%
    \noindent
    \makebox[\textwidth][c]{
        \includegraphics[width=0.9\paperwidth,keepaspectratio]{#1}%
    }%
}%


\ifwithsols
\title{حل امتحان مفمار في علم الحاسوب للصف الثامن - بايثون $2022$}
\else
\title{امتحان مفمار في علم الحاسوب للصف الثامن - بايثون $2022$}
\fi

\begin{document}

\maketitle
\thispagestyle{fancy}

\ifwithsols
\else
\begin{boxCode}
\begin{itemize}[nosep]
    \item \textbf{الفصل الأول (إجباري - $60$ علامة):} يجب الإجابة عن \textbf{جميع} الأسئلة $1$-$5$.
    \item \textbf{الفصل الثاني (اختياري - $40$ علامة):} يجب الإجابة عن \textbf{سؤالين} من بين الأسئلة $6$-$9$.
    \item \textbf{مدة الامتحان:} $120$ دقيقة.
    \item \textbf{مواد مساعدة مسموحة:} كل مادة مساعدة، باستثناء حاسبة أو حاسوب قابل للبرمجة.
    \item اقرؤوا أسئلة الامتحان بتمعّن وأجيبوا عنها بانتباه في المكان المخصّص لذلك.
\end{itemize}
\end{boxCode}
\fi

\section*{الفصل الأول — إجباري ($60$ علامة)}
في هذا الفصل خمسة أسئلة. أجيبوا عن \textbf{جميع} الأسئلة $1$-$5$.

\begin{enumerate}[label=\textbf{السؤال \arabic*:}, wide=0pt, leftmargin=*, itemsep=2.2em]

% ---------------------------------------------------------------
\item ($12$ علامة) أمامكم جدول فيه أربعة مقاطع برنامج بلغة بايثون. لكل مقطع برنامج في الجدول يجب كتابة ما هو المُخرج الناتج في العمود الفارغ.

{\renewcommand{\arraystretch}{1.4}
\begin{longtable}{|C{1.5cm}|p{9cm}|C{4cm}|}
\hline
\textbf{البند} & \centering\textbf{كود بلغة بايثون} & \textbf{المُخرج الناتج} \\
\hline
$1$. &
\begin{pycell}
\begin{english}
\begin{minted}{python}
num1 = 3
num2 = 10
num1 = num2
num2 = num1
print("num1 = ", num1, "num2 = ", num2)
\end{minted}
\end{english}
\end{pycell}
& \ifwithsols \begin{english}\small\texttt{num1 = 10 num2 = 10}\end{english} \fi \\
\hline
$2$. &
\begin{pycell}
\begin{english}
\begin{minted}{python}
number = "12345"
print("4" in number)
\end{minted}
\end{english}
\end{pycell}
& \ifwithsols \begin{english}\texttt{True}\end{english} \fi \\
\hline
$3$. &
\begin{pycell}
\begin{english}
\begin{minted}{python}
msg1 = "He" + "l" * 2 + "o"
msg2 = " world"
print(msg1 + msg2)
print(len(msg1))
\end{minted}
\end{english}
\end{pycell}
& \ifwithsols \begin{english}\texttt{Hello world}\newline\texttt{5}\end{english} \fi \\
\hline
$4$. &
\begin{pycell}
\begin{english}
\begin{minted}{python}
number = 123
print(number // 100)
print(number % 100 % 10)
print(number % 10)
\end{minted}
\end{english}
\end{pycell}
& \ifwithsols \begin{english}\texttt{1}\newline\texttt{3}\newline\texttt{3}\end{english} \fi \\
\hline
\end{longtable}}
\filbreak

% ---------------------------------------------------------------
\item ($12$ علامة) معطاة ثلاثة متغيرات $a$، $b$، $c$ وقيمها كالتالي:
$$a = 10, \quad b = 15, \quad c = 0$$
احسبوا قيمة كل واحد من التعابير البوليانية في الجدول، واكتبوا القيمة (\textenglish{True} / \textenglish{False}) في العمود الفارغ.

{\renewcommand{\arraystretch}{1.4}
\begin{longtable}{|C{1.5cm}|>{\raggedright\arraybackslash}p{9cm}|C{3.5cm}|}
\hline
\textbf{البند} & \centering\textbf{تعبير بولياني} & \textbf{قيمة التعبير} \\
\hline
$1$. & \textenglish{\texttt{a >= b or b == c}} & \ifwithsols \textenglish{False} \fi \\ \hline
$2$. & \textenglish{\texttt{a \% 5 == 0 and b // 5 > c}} & \ifwithsols \textenglish{True} \fi \\ \hline
$3$. & \textenglish{\texttt{a < 0 or (b == c + 15 and a >= c)}} & \ifwithsols \textenglish{True} \fi \\ \hline
$4$. & \textenglish{\texttt{a \% 2 == c \% 2 and b \% 3 > a \% 3}} & \ifwithsols \textenglish{False} \fi \\ \hline
$5$. & \textenglish{\texttt{a != b and b != c}} & \ifwithsols \textenglish{True} \fi \\ \hline
\end{longtable}}
\filbreak

% ---------------------------------------------------------------
\item ($12$ علامة) أمامكم جدول فيه أربعة مقاطع برنامج. لكل مقطع برنامج في الجدول يجب كتابة ما هو المُخرج الناتج في العمود الفارغ.

{\renewcommand{\arraystretch}{1.4}
\begin{longtable}{|C{1.5cm}|p{9cm}|C{4cm}|}
\hline
\textbf{البند} & \centering\textbf{كود بلغة بايثون} & \textbf{المُخرج الناتج} \\
\hline
$1$. &
\begin{pycell}
\begin{english}
\begin{minted}{python}
i = 1
while i < 4:
    print('wow')
    i += 1
\end{minted}
\end{english}
\end{pycell}
& \ifwithsols \begin{english}\texttt{wow}\newline\texttt{wow}\newline\texttt{wow}\end{english} \fi \\
\hline
$2$. &
\begin{pycell}
\begin{english}
\begin{minted}{python}
for item in range(6):
    if item % 2 == 1:
        print(item * item)
\end{minted}
\end{english}
\end{pycell}
& \ifwithsols \begin{english}\texttt{1}\newline\texttt{9}\newline\texttt{25}\end{english} \fi \\
\hline
$3$. &
\begin{pycell}
\begin{english}
\begin{minted}{python}
for item in range(9, 0, -3):
    print(item)
\end{minted}
\end{english}
\end{pycell}
& \ifwithsols \begin{english}\texttt{9}\newline\texttt{6}\newline\texttt{3}\end{english} \fi \\
\hline
$4$. &
\begin{pycell}
\begin{english}
\begin{minted}{python}
x = 5
i = 1
while i < 4:
    x -= 1
    i += 1
print(i, x)
\end{minted}
\end{english}
\end{pycell}
& \ifwithsols \begin{english}\texttt{4 2}\end{english} \fi \\
\hline
\end{longtable}}
\filbreak

% ---------------------------------------------------------------
\item ($12$ علامة) أمامكم عمودان من الحلقات. يجب الملاءمة بين أزواج الحلقات المتكافئة (بحيث ينتج في الحلقتين نفس المُخرج تمامًا).

\textbf{مثال:} الحلقتان التاليتان متكافئتان لأنّ لهما نفس المُخرج (من اليسار إلى اليمين): \textenglish{\texttt{5 6 7 8 9 10}}

{\renewcommand{\arraystretch}{1.3}
\begin{longtable}{|p{6.5cm}|p{6.5cm}|}
\hline
\begin{pycell}
\begin{english}
\begin{minted}{python}
for item in range(5, 11):
    print(item, end=" ")
\end{minted}
\end{english}
\end{pycell}
&
\begin{pycell}
\begin{english}
\begin{minted}{python}
for item in range(6):
    print(item + 5, end=" ")
\end{minted}
\end{english}
\end{pycell}
\\ \hline
\end{longtable}}

لائموا بين كل حلقة في العمود "أ" والحلقة المكافئة لها في العمود "ب":

{\renewcommand{\arraystretch}{1.3}
\begin{longtable}{|C{0.9cm}|p{5.8cm}|C{0.9cm}|p{5.8cm}|}
\hline
\multicolumn{2}{|c|}{\textbf{العمود أ}} & \multicolumn{2}{c|}{\textbf{العمود ب}} \\
\hline
$1$ &
\begin{pycell}
\begin{english}
\begin{minted}{python}
for x in range(100, 200, 2):
    print(x)
\end{minted}
\end{english}
\end{pycell}
& $a$ &
\begin{pycell}
\begin{english}
\begin{minted}{python}
x = 0
while x < 9:
    print(x)
    x += 1
\end{minted}
\end{english}
\end{pycell}
\\ \hline
$2$ &
\begin{pycell}
\begin{english}
\begin{minted}{python}
for x in range(1, 10):
    print(x - 1)
\end{minted}
\end{english}
\end{pycell}
& $b$ &
\begin{pycell}
\begin{english}
\begin{minted}{python}
for x in range(7, 8):
    print(x)
\end{minted}
\end{english}
\end{pycell}
\\ \hline
$3$ &
\begin{pycell}
\begin{english}
\begin{minted}{python}
for x in range(4, 5):
    print(x + 3)
\end{minted}
\end{english}
\end{pycell}
& $c$ &
\begin{pycell}
\begin{english}
\begin{minted}{python}
for x in range(100, 200):
    if x % 2 == 0:
        print(x)
\end{minted}
\end{english}
\end{pycell}
\\ \hline
$4$ &
\begin{pycell}
\begin{english}
\begin{minted}{python}
for x in range(-5, 6):
    print(x + 6)
\end{minted}
\end{english}
\end{pycell}
& $d$ &
\begin{pycell}
\begin{english}
\begin{minted}{python}
x = 1
while x < 12:
    print(x)
    x += 1
\end{minted}
\end{english}
\end{pycell}
\\ \hline
\end{longtable}}

\ifwithsols
\begin{boxSolution}
\begin{itemize}
\item الحلقة $1$ تلائم الحلقة $c$ (كلاهما يطبع الأعداد الزوجية من $100$ إلى $198$).
\item الحلقة $2$ تلائم الحلقة $a$ (كلاهما يطبع الأعداد من $0$ إلى $8$).
\item الحلقة $3$ تلائم الحلقة $b$ (كلاهما يطبع العدد $7$ فقط).
\item الحلقة $4$ تلائم الحلقة $d$ (كلاهما يطبع الأعداد من $1$ إلى $11$).
\end{itemize}
\end{boxSolution}
\fi
\filbreak

% ---------------------------------------------------------------
\item ($12$ علامة) أمامكم أربعة مقاطع برنامج تتضمّن عملية باسم \textenglish{calc\_avg}.
الهدف من كل واحد من مقاطع البرنامج هو استقبال علامات طلاب الصف الثامن في امتحان الرياضيات من المستخدم حتى تُستقبل علامة سالبة أو صفر، وحساب معدّل العلامات الموجبة التي استُقبلت.
(انتبهوا: العلامة السالبة أو الصفر التي تُستقبل في نهاية الإدخال ليست جزءًا من حساب المعدّل.)
في كل مقطع برنامج وقع خطأ منطقي واحد، ولذلك لا يحقّق المقطع هدفه. اكتبوا التعليمة الخاطئة واشرحوا كيف يمكن تصحيحها.
(ملاحظة: الخطأ ليس خطأ وقت تشغيل. يمكن الافتراض أنّ المُدخل سليم ويشمل أعدادًا صحيحة فقط، وتوجد علامة موجبة واحدة على الأقل في سلسلة القيم المستقبَلة.)

\begin{enumerate}[label=\textbf{مقطع \alph*:}, wide=0pt, leftmargin=*, itemsep=1.5em]

\item .
\begin{boxCode}
\begin{english}
\begin{minted}{python}
def calc_avg():
    sum = 0
    count = 0
    grade = int(input("Please enter Math's grade:"))
    while grade > 0:
        sum += grade
        count += 1
    avg = sum / count
    print(avg)
\end{minted}
\end{english}
\end{boxCode}
\ifwithsols
\begin{boxSolution}
\textbf{الخطأ:} تنقص تعليمة استقبال علامة جديدة داخل حلقة \textenglish{while}، مما يؤدّي إلى حلقة لا نهائية إذا كانت العلامة الأولى موجبة. \\
\textbf{التصحيح:} إضافة السطر \textenglish{\texttt{grade = int(input("Please enter Math's grade:"))}} في نهاية جسم حلقة \textenglish{while}.
\end{boxSolution}
\else
\textbf{ما هي التعليمة الخاطئة (أو الناقصة)؟} \underline{\hspace{8.5cm}} \\
\textbf{كيف يجب تصحيحها؟} \underline{\hspace{10cm}}
\fi
\filbreak

\item .
\begin{boxCode}
\begin{english}
\begin{minted}{python}
def calc_avg():
    sum = 0
    count = 0
    grade = int(input("Please enter Math's grade:"))
    while grade < 0:
        sum += grade
        count += 1
        grade = int(input("Please enter Math's grade:"))
    avg = sum / count
    print(avg)
\end{minted}
\end{english}
\end{boxCode}
\ifwithsols
\begin{boxSolution}
\textbf{التعليمة الخاطئة:} \textenglish{\texttt{while grade < 0:}} \\
\textbf{التصحيح:} تغيير شرط الحلقة إلى \textenglish{\texttt{while grade > 0:}} لأنّ الحلقة يجب أن تستمر طالما أنّ العلامة موجبة وتتوقف عند إدخال علامة سالبة أو صفر.
\end{boxSolution}
\else
\textbf{ما هي التعليمة الخاطئة؟} \underline{\hspace{9.5cm}} \\
\textbf{كيف يجب تصحيحها؟} \underline{\hspace{10cm}}
\fi
\filbreak

\item .
\begin{boxCode}
\begin{english}
\begin{minted}{python}
def calc_avg():
    sum = 0
    count = 0
    grade = int(input("Please enter Math's grade:"))
    while grade > 0:
        sum += grade
        count += 1
        grade = int(input("Please enter Math's grade:"))
        avg = sum / count
    print(avg)
\end{minted}
\end{english}
\end{boxCode}
\ifwithsols
\begin{boxSolution}
\textbf{التعليمة الخاطئة:} حساب المعدّل \textenglish{\texttt{avg = sum / count}} موجود داخل حلقة \textenglish{while} (إزاحة زائدة). رغم أنّه يحسب القيمة الصحيحة في التكرار الأخير، إلّا أنّ حساب المعدّل يجب أن يتمّ مرة واحدة فقط بعد انتهاء الحلقة. \\
\textbf{التصحيح:} إخراج السطر \textenglish{\texttt{avg = sum / count}} ليكون خارج حلقة \textenglish{while} (بنفس مستوى إزاحة \textenglish{while} و \textenglish{print}).
\end{boxSolution}
\else
\textbf{ما هي التعليمة الخاطئة؟} \underline{\hspace{9.5cm}} \\
\textbf{كيف يجب تصحيحها؟} \underline{\hspace{10cm}}
\fi
\filbreak

\item .
\begin{boxCode}
\begin{english}
\begin{minted}{python}
def calc_avg():
    sum = 0
    count = 0
    grade = int(input("Please enter Math's grade:"))
    while grade > 0:
        sum += grade
        count -= 1
        grade = int(input("Please enter Math's grade:"))
    avg = sum / count
    print(avg)
\end{minted}
\end{english}
\end{boxCode}
\ifwithsols
\begin{boxSolution}
\textbf{التعليمة الخاطئة:} \textenglish{\texttt{count -= 1}} (إنقاص عدّاد العلامات بـ $1$ بدل زيادته). \\
\textbf{التصحيح:} تغيير السطر إلى \textenglish{\texttt{count += 1}} لعدّ العلامات الموجبة بشكل صحيح.
\end{boxSolution}
\else
\textbf{ما هي التعليمة الخاطئة؟} \underline{\hspace{9.5cm}} \\
\textbf{كيف يجب تصحيحها؟} \underline{\hspace{10cm}}
\fi
\filbreak

\end{enumerate}
\filbreak

\end{enumerate}

\clearpage
\section*{الفصل الثاني — اختياري ($40$ علامة)}
أجيبوا عن \textbf{سؤالين} من بين الأسئلة $6$-$9$ (لكل سؤال — $20$ علامة).

\begin{enumerate}[resume, label=\textbf{السؤال \arabic*:}, wide=0pt, leftmargin=*, itemsep=2.2em]

% ---------------------------------------------------------------
\item ($20$ علامة) \textbf{العدد المثالي} هو عدد يساوي مجموع جميع قواسمه الموجبة (لا يشمل العدد نفسه).
\begin{itemize}
\item \textbf{مثال:} العدد $28$ هو عدد مثالي لأنّه يساوي مجموع جميع قواسمه: $28 = 1 + 2 + 4 + 7 + 14$.
\item القاسم هو عدد يقسم العدد المقسوم مع باقي قسمة صفر. مثلًا: $4$ هو قاسم لـ $24$، لكنّه ليس قاسمًا لـ $23$.
\end{itemize}

\begin{enumerate}[label=\alph*.]
\item حسب التعريف أعلاه، هل $6$ هو عدد مثالي؟ إذا نعم، اعرضوا تعبيرًا رياضيًا (كما هو معروض في المثال) يُظهر ذلك؛ وإلّا اشرحوا.
\ifwithsols
\begin{boxSolution}
\textbf{نعم}، العدد $6$ هو عدد مثالي لأنّ قواسمه (بدون الـ $6$ نفسه) هي $1$، $2$، $3$، ومجموعها يساوي:
$$6 = 1 + 2 + 3$$
\end{boxSolution}
\else
\\ \underline{\hspace{14cm}}
\fi
\filbreak

\item أمامكم مقطع برنامج يستقبل عددًا من المستخدم ويفحص ما إذا كان العدد مثاليًا. أكملوا الفراغات $(1)$-$(5)$ في المقطع:

\begin{boxCode}
\begin{english}
\begin{minted}{python}
def perfect_number():
    num = int(input("Please enter a number:"))
    sum = (1)______
    x = 1
    while x < num:
        if (2)______________:
            sum = sum + (3)_________
        x = (4)_________
    if (5)______________:
        print("Perfect number ")
    else:
        print("Not perfect number ")

perfect_number()
\end{minted}
\end{english}
\end{boxCode}

\ifwithsols
\begin{boxSolution}
\begin{enumerate}[label=(\arabic*)]
\item $0$
\item \textenglish{\texttt{num \% x == 0}}
\item $x$
\item $x + 1$
\item \textenglish{\texttt{sum == num}}
\end{enumerate}
\end{boxSolution}
\fi
\filbreak
\end{enumerate}
\filbreak

% ---------------------------------------------------------------
\item ($20$ علامة) مقطع البرنامج الذي أمامكم يستقبل رقم هوية مكوّنًا من $9$ منازل (\textenglish{id\_number}) ومنزلة إضافية (\textenglish{digit})، ويطبع جميع الأماكن (المواقع) في رقم الهوية التي تظهر فيها المنزلة \textenglish{digit} (من اليمين إلى اليسار، بدءًا من المكان $1$).
إذا لم تظهر المنزلة في رقم الهوية، تطبع العملية $-1$.
\begin{itemize}
\item \textbf{مثال $1$:} لرقم الهوية $334429969$ والمنزلة $9$ — تُطبع الأماكن التالية: $1$، $3$، $4$.
\item \textbf{مثال $2$:} لرقم الهوية $123457899$ والمنزلة $6$ — يُطبع $-1$ لأنّ المنزلة $6$ غير موجودة في رقم الهوية.
\end{itemize}
أكملوا الفراغات $(1)$-$(7)$ في المقطع لكي ينفّذ المطلوب:

\begin{boxCode}
\begin{english}
\begin{minted}{python}
def find_location():
    id_number = int(input("Please enter ID number:"))
    digit = int(input("Please enter a digit:"))
    found = False
    for index in range((1)_______, (2)_______):
        if (3)__________________:
            print((4)____________)
            found = (5)_____________
        id_number = (6)_______________
    if found == False:
        print((7)__________)

find_location()
\end{minted}
\end{english}
\end{boxCode}

\ifwithsols
\begin{boxSolution}
\begin{enumerate}[label=(\arabic*)]
\item $1$
\item $10$
\item \textenglish{\texttt{id\_number \% 10 == digit}}
\item \textenglish{\texttt{index}}
\item \textenglish{\texttt{True}}
\item \textenglish{\texttt{id\_number // 10}}
\item $-1$
\end{enumerate}
\end{boxSolution}
\fi
\filbreak

% ---------------------------------------------------------------
\item ($20$ علامة) أمامكم أربعة مقاطع برنامج تستعمل المكتبة \textenglish{turtle}. كل مقطع برنامج يرسم شكلًا مختلفًا. لائموا بين كل كود والشكل الذي يرسمه.
افترضوا أنّه في بداية كل كود تظهر الأسطر التالية المشتركة لجميع المقاطع:
\begin{boxCode}
\begin{english}
\begin{minted}{python}
import turtle
wn = turtle.Screen()
player = turtle.Turtle()
player.shape("turtle")
\end{minted}
\end{english}
\end{boxCode}

{\renewcommand{\arraystretch}{1.3}
\begin{longtable}{|C{0.9cm}|p{7.2cm}|C{0.9cm}|C{5.2cm}|}
\hline
\multicolumn{2}{|c|}{\textbf{كود بايثون}} & \multicolumn{2}{c|}{\textbf{الشكل}} \\
\hline
$1)$ &
\begin{pycell}
\begin{english}
\begin{minted}{python}
for x in range(360):
    player.left(1)
    player.forward(1)
turtle.mainloop()
\end{minted}
\end{english}
\end{pycell}
& $a)$ &
\begin{english}
\vspace{0.1cm}
\begin{tikzpicture}[scale=0.45]
\draw[thick] (0,0) -- (0,1) -- (1,1) -- (1,2) -- (2,2) -- (2,3) -- (3,3) -- (3,4) -- (4,4);
\filldraw (4,4) circle (3pt);
\end{tikzpicture}
\end{english}
\\ \hline
$2)$ &
\begin{pycell}
\begin{english}
\begin{minted}{python}
for index in range(3):
    player.forward(100)
    player.left(360 / 3)
turtle.mainloop()
\end{minted}
\end{english}
\end{pycell}
& $b)$ &
\begin{english}
\vspace{0.1cm}
\begin{tikzpicture}[scale=0.55]
\draw[thick] (0,0) -- (3,0) -- (1.5,2.6) -- cycle;
\filldraw (0,0) circle (3pt);
\end{tikzpicture}
\end{english}
\\ \hline
$3)$ &
\begin{pycell}
\begin{english}
\begin{minted}{python}
player.left(180)
player.forward(80)
player.left(90)
player.forward(70)
for index in range(4):
    player.forward(80)
    player.left(90)
turtle.mainloop()
\end{minted}
\end{english}
\end{pycell}
& $c)$ &
\begin{english}
\vspace{0.1cm}
\begin{tikzpicture}[scale=0.55]
\draw[thick] (0,1.5) circle (1.5cm);
\filldraw (0,0) circle (3pt);
\end{tikzpicture}
\end{english}
\\ \hline
$4)$ &
\begin{pycell}
\begin{english}
\begin{minted}{python}
for index in range(4):
    player.left(90)
    player.forward(50)
    player.right(90)
    player.forward(50)
turtle.mainloop()
\end{minted}
\end{english}
\end{pycell}
& $d)$ &
\begin{english}
\vspace{0.1cm}
\begin{tikzpicture}[scale=0.45]
\draw[thick] (2,4) -- (0,4) -- (0,2.2) -- (0,0) -- (2,0) -- (2,2.2) -- (0,2.2);
\filldraw (0,2.2) circle (3pt);
\end{tikzpicture}
\end{english}
\\ \hline
\end{longtable}}

\ifwithsols
\begin{boxSolution}
\begin{itemize}
\item الكود $1$ يلائم الشكل $c$ (دائرة، لأنّه يتحرك خطوة واحدة ويدور درجة واحدة $360$ مرة).
\item الكود $2$ يلائم الشكل $b$ (مثلث متساوي الأضلاع، لأنّه يكرر $3$ مرات وبزاوية دوران $120^\circ$).
\item الكود $3$ يلائم الشكل $d$ (يتحرك يسارًا ثم للأسفل، ثم يرسم مربّعًا).
\item الكود $4$ يلائم الشكل $a$ (درج صاعد مكوّن من $4$ درجات).
\end{itemize}
\end{boxSolution}
\fi
\filbreak

% ---------------------------------------------------------------
\item ($20$ علامة) \textbf{"نص بقفزة $2$"} هو نص يُبنى من نص موجود عن طريق تخطّي حرفين اثنين في كل مرة (أي أخذ حرف وتخطّي حرفين، وهكذا).
\begin{itemize}
\item \textbf{مثال:} للنص التالي: \textenglish{\texttt{"paryertgfhyjogqn"}}، ينتج النص بقفزة $2$ التالي: \textenglish{\texttt{"python"}}.
\end{itemize}
مقطع البرنامج الذي أمامكم يستقبل نصًا \textenglish{st} من المستخدم، ويعرض "نصًا بقفزة $2$". أكملوا الأجزاء الناقصة $(1)$-$(6)$ في المقطع:

\begin{boxCode}
\begin{english}
\begin{minted}{python}
def create_skip2str():
    (1)___ = input("Please enter a string:")
    skip2str = ""
    for index in range((2)____, (3)________, (4)____):
        skip2str = skip2str + (5)_________
    print((6)________)

create_skip2str()
\end{minted}
\end{english}
\end{boxCode}

\ifwithsols
\begin{boxSolution}
\begin{enumerate}[label=(\arabic*)]
\item \textenglish{\texttt{st}}
\item $0$
\item \textenglish{\texttt{len(st)}}
\item $3$ (لأنّ تخطّي حرفين يعني القفز بخطوة مقدارها $3$ مؤشرات: $0, 3, 6, \dots$)
\item \textenglish{\texttt{st[index]}}
\item \textenglish{\texttt{skip2str}}
\end{enumerate}
\end{boxSolution}
\fi
\filbreak

\end{enumerate}

\vspace{1cm}
\begin{flushleft}
أرجو لكم وقتًا ممتعًا.

الأستاذ محمود اغبارية.
\end{flushleft}

\end{document}
